% !TEX TS-program = lualatex
\section{Постановка задачи}

Целью работы является разработка интеллектуального мобильного фитнес-приложения для операционной системы Android, предназначенного для организации тренировочного процесса, формирования персонализированных программ занятий, отслеживания физической активности пользователя и анализа факторов восстановления организма.

Для достижения поставленной цели необходимо решить следующие задачи:

\begin{enumerate}
\item Выполнить проектирование и реализацию пользовательского интерфейса приложения, обеспечивающего удобное взаимодействие с функциональными возможностями системы, навигацию между экранами, просмотр тренировок и упражнений, а также создание индивидуальных тренировочных программ.

\item Разработать подсистему хранения и обработки данных, включающую проектирование локальной базы данных для хранения информации о пользователях, тренировках, упражнениях и параметрах приложения, а также механизмы обмена данными с серверной частью системы.

\item Реализовать бизнес-логику приложения, обеспечивающую проведение тренировок, учёт результатов занятий, мониторинг показателей сна и формирование рекомендаций на основе собранных данных.

\item Спроектировать и разработать серверную часть приложения, предоставляющую программный интерфейс взаимодействия (API), реализующую механизмы аутентификации и авторизации пользователей, а также безопасное хранение и синхронизацию пользовательских данных.

% \item Провести тестирование разработанного программного обеспечения и оценить корректность функционирования его основных компонентов.

\end{enumerate}
